% ============================================================
% Capítulo 5 — Configuración
% ============================================================
\chapter{Configuración del proyecto}
\label{cap:configuracion}

Todo el comportamiento del proyecto está controlado por
\file{config.py}. Este archivo define un diccionario de
configuración (\texttt{DotDict}) que se pasa a cada módulo y
permite modificar el experimento sin tocar el código.

\section{La clase \texttt{DotDict}}

\begin{minted}[numbers, firstnumber=14, linerange=14-17]{python}
class DotDict(dict):
    __getattr__ = dict.get
    __setattr__ = dict.__setitem__
    __delattr__ = dict.__delitem__
\end{minted}

\noindent Es un \texttt{dict} con acceso por atributo:
\texttt{configs.lr} es equivalente a \texttt{configs['lr']}.
Esto permite usar tanto \texttt{config.feature\_length} como
\texttt{config['feature\_length']} en el resto del código.

\section{Función principal: \texttt{sys\_configuration}}

Es la única función exportada. Acepta tres parámetros y devuelve
un \texttt{DotDict}:

\begin{minted}[numbers, firstnumber=20, linerange=20-31]{python}
def sys_configuration(platform_name: str = platform.node(),
                     dataset_name: str = "CUHK-PEDES",
                     dataset_source: str = "huggingface",
                     show_configs: bool = False) -> DotDict:
    """
    Args.:
        platform_name: hostname of the machine
        dataset_name:  "CUHK-PEDES" (only one supported for now)
        dataset_source: "huggingface" or "original"
    """
\end{minted}

\noindent Los argumentos se pueden sobreescribir desde CLI en
\texttt{train.py} (\fline{train.py}{56}).

\section{Tabla maestra de parámetros}

\subsection*{Modelo de lenguaje}

\begin{longtable}{p{3.5cm}p{2.5cm}p{8cm}}
\toprule
\textbf{Parámetro} & \textbf{Valor} & \textbf{Efecto / ubicación} \\
\midrule
\endhead
\texttt{dropout\_rate}
  & 0.30
  & Dropout tras el embedding textual. \\
  & & \fline{config.py}{37}, \fline{model/language_network.py}{23}. \\
\midrule
\texttt{num\_layers}
  & 1
  & Número de capas del \gru{} bidireccional. \\
  & & \fline{config.py}{38}, \fline{model/language_network.py}{24}. \\
\midrule
\texttt{tokens\_length\_max}
  & 100
  & Longitud fija (con padding) de cada caption tokenizada. \\
  & & \fline{config.py}{39}, \fline{datasets/bases.py}{96}. \\
\midrule
\texttt{feature\_length}
  & 1024
  & Dimensión del espacio de embedding conjunto. \\
  & & \fline{config.py}{40}, \fline{model/textreidnet.py}{43}. \\
\midrule
\texttt{tokenizer\_type}
  & 'bert'
  & Tipo de tokenizer. Otros: 'simple\_tokenizer', 'tiktoken\_*'. \\
  & & \fline{config.py}{41}, \fline{datasets/bases.py}{97}. \\
\midrule
\texttt{vocab\_size}
  & 29\,610
  & 29\,609 (BERT) + 1 (padding). Padding idx = 0. \\
  & & \fline{config.py}{42}, \fline{model/language_network.py}{22}. \\
\midrule
\texttt{embedding\_dim}
  & 512
  & Dimensión del embedding textual. \\
  & & \fline{config.py}{43}, \fline{model/language_network.py}{22}. \\
\bottomrule
\caption{Parámetros del modelo de lenguaje.}
\label{tab:cfg-lang}
\end{longtable}

\subsection*{Modelo visual}

\begin{longtable}{p{3.5cm}p{2.5cm}p{8cm}}
\toprule
\textbf{Parámetro} & \textbf{Valor} & \textbf{Efecto / ubicación} \\
\midrule
\endhead
\texttt{CUHKPEDES\_image\_size}
  & (384, 128)
  & Tamaño \emph{declarado} en el paper. \textbf{No se usa}
        en el dataloader actual. \\
  & & \fline{config.py}{48}. \\
\midrule
\texttt{MHPV2\_image\_size}
  & (512, 512)
  & Tamaño del canvas en el que se pegan las imágenes. \\
  & & \fline{config.py}{49}, \fline{datasets/bases.py}{58}. \\
\midrule
\texttt{MHPV2\_means}
  & [0.355, 0.330, 0.311]
  & Medias para \texttt{T.Normalize} en el dataloader. \\
  & & \fline{config.py}{50}. \\
\midrule
\texttt{MHPV2\_stds}
  & [0.348, 0.334, 0.330]
  & Desviaciones para \texttt{T.Normalize}. \\
  & & \fline{config.py}{51}. \\
\midrule
\texttt{CUHKPEDES\_means}
  & [0.485, 0.456, 0.406]
  & Medias de ImageNet (no usadas en este proyecto). \\
  & & \fline{config.py}{52}. \\
\midrule
\texttt{CUHKPEDES\_stds}
  & [0.229, 0.224, 0.225]
  & Desviaciones de ImageNet (no usadas). \\
  & & \fline{config.py}{53}. \\
\bottomrule
\caption{Parámetros del modelo visual.}
\label{tab:cfg-vis}
\end{longtable}

\subsection*{Entrenamiento}

\begin{longtable}{p{3.5cm}p{2.5cm}p{8cm}}
\toprule
\textbf{Parámetro} & \textbf{Valor} & \textbf{Efecto / ubicación} \\
\midrule
\endhead
\texttt{epoch}
  & 60
  & Número de epochs total. \\
  & & \fline{config.py}{71}, \fline{train.py}{113}. \\
\midrule
\texttt{adam\_alpha}
  & 0.90
  & $\beta_1$ de AdamW. \\
  & & \fline{config.py}{72}, \fline{train.py}{98}. \\
\midrule
\texttt{adam\_beta}
  & 0.999
  & $\beta_2$ de AdamW. \\
  & & \fline{config.py}{73}, \fline{train.py}{98}. \\
\midrule
\texttt{epoch\_decay}
  & [20, 40]
  & Épocas donde el LR scheduler decae el learning rate. \\
  & & \fline{config.py}{74}, \fline{train.py}{99}. \\
\midrule
\texttt{lr}
  & 0.001
  & Learning rate inicial. \\
  & & \fline{config.py}{75}, \fline{train.py}{98}. \\
\midrule
\texttt{patience}
  & 3
  & (No usado actualmente; sería early stopping). \\
  & & \fline{config.py}{76}. \\
\midrule
\texttt{val\_dataset}
  & 'test'
  & Split usado durante validación. \\
  & & \fline{config.py}{77}. \\
\midrule
\texttt{seed}
  & 3407
  & Semilla para reproducibilidad. \\
  & & \fline{config.py}{79}, \fline{train.py}{71}. \\
\midrule
\texttt{batch\_size}
  & 8 (default)
  & Tamaño de batch (4 en Jetson). \\
  & & \fline{config.py}{123,126,129}. \\
\midrule
\texttt{num\_workers}
  & 4 (default)
  & Hilos del DataLoader. \\
  & & \fline{config.py}{122,125,128}. \\
\bottomrule
\caption{Parámetros de entrenamiento.}
\label{tab:cfg-train}
\end{longtable}

\subsection*{Pérdidas}

\begin{longtable}{p{3.5cm}p{2.5cm}p{8cm}}
\toprule
\textbf{Parámetro} & \textbf{Valor} & \textbf{Efecto / ubicación} \\
\midrule
\endhead
\texttt{margin}
  & 0.5
  & Margen del triplet-style RankingLoss. \\
  & & \fline{config.py}{86}, \fline{evaluation/ranking_loss.py}{41}. \\
\midrule
\texttt{ranking\_loss\_alpha}
  & 1.0
  & Peso $\alpha$ de la RankingLoss en la pérdida total. \\
  & & \fline{config.py}{87}, \fline{train.py}{145}. \\
\midrule
\texttt{identity\_loss\_beta}
  & 1.0
  & Peso $\beta$ de la IdentityLoss en la pérdida total. \\
  & & \fline{config.py}{88}, \fline{train.py}{146}. \\
\bottomrule
\caption{Parámetros de funciones de pérdida.}
\label{tab:cfg-loss}
\end{longtable}

\subsection*{Paths y logging}

\begin{longtable}{p{3.5cm}p{5cm}p{5.5cm}}
\toprule
\textbf{Parámetro} & \textbf{Valor} & \textbf{Ubicación} \\
\midrule
\endhead
\texttt{train\_log\_path}
  & \file{logs/train.log}
  & \fline{config.py}{93} \\
\texttt{test\_log\_path}
  & \file{logs/test.log}
  & \fline{config.py}{94} \\
\texttt{model\_save\_path}
  & \file{data/checkpoints}
  & \fline{config.py}{96} \\
\texttt{plot\_save\_path}
  & \file{data/plots}
  & \fline{config.py}{97} \\
\texttt{dataset\_path} (HF)
  & \file{data/CUHK-PEDES-HF}
  & \fline{config.py}{108} \\
\texttt{dataset\_path} (orig)
  & \file{data/CUHK-PEDES}
  & \fline{config.py}{112} \\
\bottomrule
\caption{Paths del proyecto.}
\label{tab:cfg-paths}
\end{longtable}

\section{Override por línea de comandos}

\file{train.py} permite sobreescribir los parámetros más
importantes:

\begin{minted}[basicstyle=\ttfamily\small, frame=single]{python}
python train.py \
    --dataset_source huggingface \
    --epochs 60 \
    --batch_size 8 \
    --lr 0.001 \
    --seed 3407 \
    --output_dir ./data/checkpoints
\end{minted}

\noindent Cada flag corresponde a un campo del config que se
sobreescribe justo después de llamar a \texttt{sys\_configuration}
(\fline{train.py}{56-69}).

\section{Configuración por plataforma}

Hay una rama de detección automática de plataforma basada en el
hostname (\fline{config.py}{121-129}):

\begin{minted}[numbers, firstnumber=121, linerange=121-129]{python}
if platform_name in ('PC-SIM', 'deeplearning', 'ultron', 'yrsn'):
    configs['num_workers'] = 4
    configs['batch_size'] = 8
elif platform_name == 'nano':
    configs['num_workers'] = 2
    configs['batch_size'] = 4
else:
    configs['num_workers'] = 4
    configs['batch_size'] = 8
\end{minted}

\noindent Si tu hostname es ``nano'' (Jetson Nano), el batch se
reduce automáticamente a 4. Para cualquier otro caso, usa 8.

\section{Reproducibilidad}

\texttt{set\_seed}
(\fline{utils/miscellaneous_utils.py}{18-30}) fija:

\begin{itemize}
  \item \texttt{PYTHONHASHSEED}.
  \item \texttt{random.seed}, \texttt{numpy.random.seed}.
  \item \texttt{torch.manual\_seed} y \texttt{torch.cuda.manual\_seed}.
  \item \texttt{cudnn.deterministic=True}, \texttt{cudnn.benchmark=False}.
\end{itemize}

\noindent Esto garantiza resultados idénticos entre ejecuciones
en la misma máquina. \textbf{No} garantiza resultados idénticos
entre máquinas distintas.

\section{Ejemplo de inspección}

Para imprimir la configuración completa desde un script:

\begin{minted}[basicstyle=\ttfamily\small, frame=single]{python}
from config import sys_configuration
cfg = sys_configuration(dataset_source='huggingface')
print(cfg.feature_length)    # 1024
print(cfg.batch_size)        # 8
print(cfg.vocab_size)        # 29610
\end{minted}

\noindent Esta es la mejor forma de verificar que la
configuración es la esperada antes de entrenar.
